\section{Preserving compact-sublevel poles in every power class}
\label{sec:r11-logarithmic-pole-representatives}

Keep the arbitrary original background, the strict analytic
Laurent box $V$, its closed box $C$, and its barrier $\rho$
from (TC1)--(TC4). In particular $\overline V$ is contained
in the interior of one original affinoid, and $S$ extends to
an ordinary neighborhood of $C$. Put $q=n-p\geq1$.
All assertions below concern the actual given representatives;
they change a potential by an explicit formula, rather than
changing it on an exceptional set.

\begin{lemma}[Logarithmic attenuation and exponential capacity tails]
\label{lem:r11-logarithmic-attenuation}
Let $u:V\to[-\infty,0]$ be extended classical psh with
$\mathrm{BG}(u,0)$, without any no-loss assumption. Set
\[
                         v=-\log(1-u),
 \tag{LP1}
\]
with the value $-\infty$ at each actual pole of $u$.
Then $v$ is negative extended classical psh, has the same
actual pole set, satisfies $\mathrm{BG}(v,0)$, and belongs
to $\mathcal F_{S,\partial}(V)$. For some finite $C_0,t_0$
depending on the stated data, one has
\[
               C_S(\{v<-t\};V)\leq C_0e^{-qt}
                                     \qquad(t\geq t_0).
 \tag{LP2}
\]
For every compact $H\Subset V$ and every $0<c<q$,
\[
                \int_H e^{-cv}\,d\mu_v<\infty.
 \tag{LP3}
\]
The integral omits the $\mu_v$-null pole set. No global
finite mass at the ordinary boundary is inferred at this stage.
\end{lemma}

\noindent\textbf{Status: \texttt{NEW\_PROVED}} for this
explicit exponential-tail attenuation, using (BQ1) as input.

\begin{proof}
The scalar map $F(r)=-\log(1-r)$ on $r\leq0$ has
$F'(r)=(1-r)^{-1}>0$ and $F''(r)=(1-r)^{-2}>0$.
Apply frozen convex increasing composition to the bounded
canonical truncations of $u$ and pass to their decreasing
limit. This proves the asserted classical property and actual
values. The identity
\[
                    \{v<-t\}=\{u<-(e^t-1)\}
 \tag{LP4}
\]
proves the compact-sublevel condition, the equality of poles,
and, by (BQ1), the bound $C(e^t-1)^{-q}$.
For sufficiently large $t$ this is (LP2).
The capacity-tail criterion (BC3)--(BC5) now puts $v$ in
the no-loss boundary class.

To prove (LP3), partition the sufficiently deep finite-value
locus into the Borel strips
\[
                     A_m=\{m\leq-v<m+1\},
                           \qquad m\in\mathbf N.
 \tag{LP5}
\]
On $A_m$, strict maximum locality identifies $\mu_v$ with
the bounded measure of $v\vee(-m-2)$. Its normalized capacity
bound gives
\[
 \mu_v(A_m)
    \leq(m+2)^q C_S(A_m;V)
    \leq(m+2)^q C_S(\{v<-m+1\};V)
    \leq C_1(m+2)^q e^{-qm}.
 \tag{LP6}
\]
The harmless shift by one keeps the last carrier strictly
open and includes any equality-level mass. Multiplication by
$e^{c(m+1)}$ and summation converge if $c<q$.
On the remaining bounded-value portion of the compact $H$,
local finiteness of $\mu_v$ completes the proof.
\end{proof}

\begin{theorem}[One actual pole representative in every finite power class]
\label{thm:r11-all-power-pole-representative}
Let $u$ satisfy Lemma~\ref{lem:r11-logarithmic-attenuation}.
For every prescribed compact $H\Subset V$ and every
$\varepsilon>0$ there exists a negative extended classical
psh function $g$ on $V$ such that
\[
 \begin{gathered}
  g=-\log(1-u)-\varepsilon
       \quad\hbox{on an ordinary neighborhood of }H
                      \hbox{ and of the pole set},\\
  \{g=-\infty\}=\{u=-\infty\},\qquad
  \mathrm{BG}(g,0),\qquad \mu_g(V)<\infty,\\
  g\in\bigcap_{\sigma>0}\mathcal E^\sigma_{S,\partial}(V).
 \end{gathered}
 \tag{LP7}
\]
More precisely, $g=A\rho$ on an ordinary boundary collar,
for a finite $A>0$, and there are finite constants such that
for all sufficiently large real $t$ and integers $j$,
\[
 \begin{split}
 C_S(\{g<-t\};V)&\leq C_2e^{-qt},\\
 \mu_g\{g\leq-t\}&\leq C_3(1+t)^q e^{-qt},\\
 \|\mu_{g,j}-\mu_g\|_{\mathrm{TV}(V)}
                  &\leq C_4(1+j)^q e^{-qj}.
 \end{split}
 \tag{LP8}
\]
For every $0<c<q$, including the actual canonical truncations,
\[
       \int_V e^{-cg}\,d\mu_g<\infty,
       \qquad
       \sup_{j\geq1}\int_Ve^{-cg_j}\,d\mu_{g,j}<\infty.
 \tag{LP9}
\]
In fact these truncation exponential moments increase to the
displayed limiting moment. The function $g$ itself is chosen
using only $u,H,\varepsilon,V,\rho$: it is independent of the
background. The same $g$ therefore satisfies (LP7)--(LP9)
simultaneously for every original symmetric middle-positive
separately closed background of every degree $p<n$ extending
to a neighborhood of $C$, with the corresponding $q=n-p$.
Constants may depend on $S,u,V,H,\varepsilon,c$; no uniformity
over backgrounds or prescribed potentials is asserted.
\end{theorem}

\noindent\textbf{Status: \texttt{NEW\_PROVED}} with the
explicit strict-box and background-extension hypotheses.

\begin{proof}
Write $v=-\log(1-u)$ and choose $M>0$. The compact set
\[
                      K=H\cup\overline{\{v<-M\}}
                              \Subset V
\]
contains the pole set. Choose an ordinary relatively compact
neighborhood $W$ of $K$ inside $V$. On the compact
$\overline W$, continuity and strict negativity of $\rho$
give an $A>0$ so large that
$A\rho<-M-\varepsilon-1$ everywhere there.
Define $g$ on the following two ordinary open sets:
\[
 \begin{cases}
 g=v-\varepsilon&\text{on }W,\\
 g=\max(v-\varepsilon,A\rho)&\text{on }V\setminus K.
 \end{cases}
 \tag{LP10}
\]
On their overlap $W\setminus K$ one has $v\geq-M$,
so the first branch strictly exceeds $A\rho$.
The definitions therefore agree on the entire ordinary overlap.
The local classical gluing property, equivalently the frozen
curvewise definition after every complete extension, proves
that $g$ is extended classical psh on $V$.
It has the same actual poles as $v$, since those lie in $K$
and the function is bounded below on $V\setminus K$.
It also gives the required equality near $H$ and all poles.

Near the ordinary boundary, the condition
$\rho>-\varepsilon/(2A)$ makes the second branch in (LP10)
strictly larger than $v-\varepsilon\leq-\varepsilon$.
Thus $g=A\rho$ on this boundary collar. Its strict negative
sublevels are consequently contained in a compact subset of
$V$, the union of $\overline W$ and an appropriate compact
negative sublevel of $\rho$.
For $t>M+\varepsilon$, the portion outside $K$ cannot belong
to $\{g<-t\}$, and the exact deep-sublevel identity is
\[
                    \{g<-t\}=\{v<-(t-\varepsilon)\}.
 \tag{LP11}
\]
Equation (LP2) gives the first estimate in (LP8).
The capacity-tail criterion therefore puts $g$ in
$\mathcal F_{S,\partial}(V)$.

Choose a compact subset of $V$ outside which $g=A\rho$.
On that compact subset $\mu_g$ is finite by the no-loss
class theorem. Outside it ordinary locality identifies the
measure with $A^q\mu_\rho$, whose total mass is finite by
(TC4). Thus $\mu_g(V)<\infty$.
In particular this argument does not assert local boundedness
of $g$ near its interior poles or attempt to apply bounded
CLN across them.

Repeat the integer-strip estimate (LP6), now with $g$ and
the global finite mass on its bounded-value part. Summing
from $m=\lfloor t\rfloor$ gives the second estimate in
(LP8): for the decreasing exponential series,
\[
 \sum_{m\geq N}(m+2)^q e^{-qm}
                         \leq C_q(N+2)^q e^{-qN}.
\]
The same strip estimate with $e^{c(m+1)}$ proves the first
assertion of (LP9) for every $c<q$. All finite powers are
bounded by a constant times this exponential, so (LP7)
follows from the definition of the power classes.

For the uncut canonical product, capacity normalization gives
\[
 \nu_{g,j}(V)
       \leq j^q C_S(\{g\leq-j\};V)
       \leq j^q C_S(\{g<-j+1\};V)
       \leq C_5j^q e^{-qj}.
 \tag{LP12}
\]
The current $\mu_{g,j}$ agrees with $\mu_g$ on $\{g>-j\}$.
The variation on the complement is at most the sum of the
discarded mass (LP12) and the second bound in (LP8). This
proves the third bound there. For any fixed $j$, the bounded
function $g_j$ equals $A\rho$ on a sufficiently narrow
ordinary collar and extends as that bounded classical function
across $\partial V$. Bounded CLN and the extension of $S$
give $\mu_{g,j}(V)<\infty$.

Finally split the weighted truncation integral at $\{g>-j\}$:
\[
 \int_V e^{-cg_j}\,d\mu_{g,j}
    =\int_{\{g>-j\}}e^{-cg}\,d\mu_g
                            +e^{cj}\nu_{g,j}(V).
 \tag{LP13}
\]
The first term is bounded by the established finite integral.
The second is $O(j^q e^{-(q-c)j})$ for large $j$, and the
finitely many remaining terms are finite. This proves (LP9).
There is also an exact distributional strengthening. Since
$g$ has no loss, (NS7) at level $-j$ gives
$\nu_{g,j}(V)=\mu_g\{g\leq-j\}$. Consequently
\[
 (-g_j)_*\mu_{g,j}
       =(y\mapsto\min(y,j))_*\bigl((-g)_*\mu_g\bigr).
 \tag{LP18}
\]
Indeed the two sides agree on the finite-value carrier $g>-j$,
and both remaining parts have the same mass at the scalar value
$j$. The omitted pole is $\mu_g$-null. Apply (LP18) to the
increasing weights $e^{c\min(y,j)}$ and use monotone convergence
to obtain the claimed monotone convergence of exponential
moments. The construction (LP10) contains no background datum,
so repeating the proof for each allowed $S$ proves the simultaneous
quantifier stated above.
Every locality or capacity normalization used the full original
tester class on the same $V$.
\end{proof}

\begin{corollary}[An exact family of complete relative pole sets]
\label{cor:r11-all-power-complete-poles}
On the stated domain, the following two families of subsets
are equal:
\[
 \begin{split}
 &\bigl\{\{u=-\infty\}:u\leq0\text{ extended classical psh},
                                  \ \mathrm{BG}(u,0)\bigr\},\\
 &\bigl\{\{g=-\infty\}:g\in\bigcap_{\sigma>0}
       \mathcal E^\sigma_{S,\partial}(V),\ \mu_g(V)<\infty\bigr\}.
 \end{split}
 \tag{LP14}
\]
Each such set, and every subset of it, has zero ordinary-open
outer $S$-capacity. The same set is avoided by every original
coefficient of $S$ and every finite product of power-class
potentials, and by every graded energy current already defined.
\end{corollary}

\noindent\textbf{Status: \texttt{NEW\_PROVED}} for the exact
pole-family equality. The nullity assertions have status
\texttt{DIRECT\_COROLLARY}; their inputs are (BQ2), (DG3), and (TC9).

\begin{proof}
Theorem~\ref{thm:r11-all-power-pole-representative} proves the
first inclusion in (LP14); the reverse is part of the printed
definition of every power class. Outer nullity follows from
(BQ2), and is inherited by arbitrary subsets. For the complete
pole set, which is Borel, (DG3) and
Corollary~\ref{cor:r11-power-class-null-ideal} give all the
stated coefficient-null assertions. Their positive variation
measures then also vanish on all measurable subsets.
\end{proof}

\begin{corollary}[Exponential coefficient and energy calculus for the representative]
\label{cor:r11-logarithmic-all-degree-exponential}
Let $g$ be the function constructed in (LP10), let
$N_k=N_k(g)$, and let $T_{k,j}=T_k(g_j)$.
For every original compact smooth complementary test $\eta$,
$0\leq k\leq q$, and $0<c<q$, the coefficient multiplier
$e^{-cg}N_k$ is locally defined and
\[
 \bigl\|\mu_{e^{-cg_j}T_{k,j}-e^{-cg}N_k,\eta}\bigr\|_{\rm TV}
       \leq C_{\eta,c}(1+j)^k e^{-(q-c)j}
 \tag{LP15}
\]
for all sufficiently large $j$. For $0\leq k\leq q-1$,
the analogous statement for the energy currents is
\[
 \bigl\|\mu_{e^{-cg_j}\Gamma_{T_{k,j}}(g_j)
                     -e^{-cg}\Gamma_{N_k}(g),\eta}\bigr\|_{\rm TV}
       \leq C'_{\eta,c}(1+j)^{k+2}e^{-(q-c)j}.
 \tag{LP16}
\]
Here the test degree in (LP16) is complementary to the energy
degree. In particular, throughout $0\leq k\leq q-1$,
\[
 D(e^{-cg}N_k)
      =-c e^{-cg}N_{k+1}
                          +c^2e^{-cg}\Gamma_{N_k}(g)
 \tag{LP17}
\]
on all original compact smooth tests. The energy currents need
not be separately closed. The signed right side of (LP17) is
separately closed as a whole by that identity, but is not
asserted to be positive. The $N_k$ themselves are the already established
separately closed middle-positive currents, with $N_0=S$.
\end{corollary}

\noindent\textbf{Status: \texttt{NEW\_PROVED}} for the
exponential weights; existence, closedness and the finite-power
energy construction are inputs from the all-degree theorems.

\begin{proof}
Since $g$ belongs to every positive power class, the original
all-degree products, energy currents and their coherent finite
level restrictions exist by
Theorems~\ref{thm:r11-power-full-recursion} and
\ref{thm:r11-power-graded-energy}. No exponential assertion
is used for their initial construction.

Fix a positive original trace control for $\eta$ supplied by
(DG2). On the strip $A_m=\{m\leq-g<m+1\}$, strict
bounded locality allows every slot to be replaced by its
truncation at $-m-2$. Equations (DG3) and (LP8) bound the
resulting trace by
\[
                         C_\eta(m+2)^k e^{-qm}.
\]
The same statement for the energy trace follows from (PG8),
with $k$ replaced by $k+2$. Energy locality here is the
already proved shifted-square/four-branch locality, not an
inference from equality of arbitrary Borel representatives.
Multiplication by $e^{c(m+1)}$ and summation prove local
exponential integrability and the limiting tail bounds.
The pole set gives zero trace mass, including at $k=0$,
by (DG3) and its outer capacity nullity.

On $\{g\leq-j\}$, the entire uncut curvature trace is
bounded by $C_\eta j^k C_S\{g\leq-j\}$, and its weighted
version by $e^{cj}$ times that quantity. The corresponding
energy bound is $C_\eta j^{k+2}e^{cj}C_S\{g\leq-j\}$.
Using the strict enlargement in (LP12) bounds both by the
displayed exponential rates. On $\{g>-j\}$, the respective
weighted measures agree exactly by the established locality.
Adding the uncut and limiting tails proves (LP15)--(LP16),
first for positive traces and then for every original coefficient.

The bounded smooth scalar chain identity is
\[
 D(e^{-cg_j}T_{k,j})
     =-c e^{-cg_j}T_{k+1,j}
                  +c^2e^{-cg_j}\Gamma_{T_{k,j}}(g_j).
\]
Pass the three terms through (LP15)--(LP16) and use
distributional continuity. This proves (LP17).
No differentiation of a restricted current is performed,
and the construction and all bounds remain valid after
each actual moment refinement of a fixed test.
\end{proof}

\paragraph{Exact advance and remaining polar gate.}
The frozen R2 trace-null witness theorem supplies a
measure-dependent attenuation for a specified nonconcentrating
trace. This section instead starts with the explicit compact
sublevel hypothesis, uses the new arbitrary-$S$ capacity tail,
and constructs one actual representative in every finite power
class, with a global boundary collar and quantitative exponential
control. It does not audit or replace that frozen theorem.

The unrestricted converse that every outer-$S$-capacity-null set
lies in a pole set in (LP14) is \texttt{FALSE} without a
nondegeneracy qualification. For example, on an original open
unit disc take $S=0$ and $E=V$. Its relative capacity is zero.
A potential with $E\subset\{u=-\infty\}$ would be identically
$-\infty$ and its strict negative sublevel would be all of $V$,
whose closure meets the ordinary boundary; it cannot satisfy
$\mathrm{BG}(u,0)$. Frozen proper-support examples give the
same obstruction for nonzero backgrounds away from their support.
An appropriate converse with explicit local nondegeneracy and
domain hypotheses remains \texttt{OPEN}. Nor is a general
ambient pluripolar set asserted to admit a global compact-sublevel
defining potential. The frozen relative-support counterexamples
continue to prevent an unrestricted identification of relative
nullity with ambient negligible or pluripolar sets.
The strict-box and extension hypotheses are used to control
the ordinary boundary; removing them is not part of this proof.
